\subsection{Structural Syntax and Function Object Creation}

A Python function is a runtime object. The \texttt{def} statement creates it, attaches a \texttt{PyCodeObject}, stores metadata, and binds a name. Calling the function later creates a new execution frame; the body does not run merely because the function was defined.

\subsubsection{The Anatomy of Function Code Objects: Inspecting Bytecode Attributes (\texttt{\_\_code\_\_}, \texttt{co\_code}, \texttt{co\_varnames}, \texttt{co\_consts}, \texttt{co\_names})}

A function object references a code object. The function is the callable wrapper; the code object holds compiled execution data: local names, constants, external names, and raw bytecode.

\begin{codebox}[]{python}{The Anatomy of Function Code Objects: Inspecting Bytecode Attributes (\textbackslash{}\_\textbackslash{}\_cod...}{str02}
def make_message(name, number=42):
    text = f"{name} says {number}"
    return text.upper()

code_obj = make_message.__code__

print(f"co_argcount:    {code_obj.co_argcount}")
print(f"co_varnames:    {code_obj.co_varnames}")
print(f"co_consts:      {code_obj.co_consts}")
print(f"co_names:       {code_obj.co_names}")
print(f"co_code length: {len(code_obj.co_code)}")
\end{codebox}


\begin{iobox}{Expected Output}
\texttt{co\_argcount:    2} \\
\texttt{co\_varnames:    ('name', 'number', 'text')} \\
\texttt{co\_consts:      (None, ' says ')} \\
\texttt{co\_names:       ('upper',)} \\
\texttt{co\_code length: 58} \\
\end{iobox}




The stable separation is: function object (callable metadata) $\rightarrow$ code object (compiled description) $\rightarrow$ \texttt{co\_code} bytes (VM instruction stream). Exact byte sequences vary by CPython version; the object chain does not.

\begin{verbatim}
import dis

def make_message(name, number=42):
    text = f"{name} says {number}"
    return text.upper()

dis.dis(make_message)
\end{verbatim}

Possible output excerpt:

\begin{verbatim}
  4           2 LOAD_FAST                0 (name)
              ...
             14 STORE_FAST               2 (text)
  5          16 LOAD_FAST                2 (text)
             18 LOAD_METHOD              0 (upper)
             40 CALL                     0
             50 RETURN_VALUE
\end{verbatim}

When the function is called, a frame executes these instructions. A function is not a label on source text; it is a runtime object connected to compiled execution data.

\subsubsection{Anonymous Expressions: The Structural Equivalence of \texttt{lambda} Functions}

A \texttt{lambda} expression creates the same kind of function object as \texttt{def}. The compiler emits a \texttt{PyCodeObject} with identical structural fields: \texttt{co\_varnames}, \texttt{co\_consts}, \texttt{co\_code}, and call behavior via \texttt{\_\_call\_\_}. Lambdas are syntactic sugar restricted to a single expression; they carry the same allocation and execution profile as named functions.

\begin{codebox}[]{python}{Anonymous Expressions: The Structural Equivalence of lambda Functions}{str01}
def add_def(x):
    return x + 42

add_lambda = lambda x: x + 42

for label, fn in [("add_def", add_def), ("add_lambda", add_lambda)]:
    print(f"{label}: type={type(fn)}, vars={fn.__code__.co_varnames}, "
          f"consts={fn.__code__.co_consts}, result={fn(8)}")
\end{codebox}


\begin{iobox}{Expected Output}
\texttt{add\_def: type=<class 'function'>, vars=('x',), consts=(None, 42), result=50} \\
\texttt{add\_lambda: type=<class 'function'>, vars=('x',), consts=(None, 42), result=50} \\
\end{iobox}




The only runtime differences are cosmetic: \texttt{\_\_name\_\_} is \texttt{'<lambda>'} and there is no docstring slot. Performance, frame allocation, and bytecode structure are equivalent. Use \texttt{lambda} only when an expression-level function is clearer; use \texttt{def} for multi-statement bodies.
